\section{The Integers}

The set $\NN$ of natural numbers with the operation of addition is not a group
because the elements do not have inverses in $\NN$. To address this, we aim to
define a larger set that also contains the inverses of natural numbers, allowing
us to define the operation of subtraction.

In this context, it makes sense to represent the pair $(m,n)$ using the arrow
$m\mapsto n$ because the difference $n-m$ represents the quantity that, when
added to $m$, gives $n$. This can be thought of as an arrow (or vector) that
starts at $m$ and ends at $n$ in the set of natural numbers. However, we must
identify pairs that represent the same difference, thus we define (see
definition~\ref{def:insieme_quoziente} of quotient set)
\[
  \ZZ \defeq (\NN \times \NN)/\sim
  \qquad\text{where}\qquad
  (m\mapsto n)\sim (m' \mapsto n')\iff n+m' = n'+m.
\]
An integer $k\in \ZZ$ will therefore be an equivalence class: $k = \Enclose{m
\mapsto n}_\sim$.
\mynote{If $n\ge m$, the equivalence class $[m\mapsto n]_\sim$ coincides with the graph of the function $\NN\colon \NN$ that translates the natural numbers by $n-m$ units to the right. Thus, it coincides with the function that adds $n-m$ to each natural number. According to how we have defined the operation of addition, it follows that $[m\mapsto n]_\sim = +(n-m)$. If instead $n<m$, the equivalence class $[m\mapsto n]_\sim$ coincides with the graph of the function that translates the natural numbers by $m-n$ units to the left and is defined only on naturals greater than or equal to $m-n$. This function is precisely the operation of subtraction, thus formally $[m\mapsto n]_\sim = -(m-n)$. Formally, we can therefore write 
\[
 \ZZ = \ENCLOSE{+n\colon n\in \NN} \cup \ENCLOSE{-n\colon n\in \NN}
\]
which perfectly aligns with our intuition.
}
From now on, we simply write $k=\Enclose{m\mapsto n}$ omitting the symbol $\sim$
that indicates equivalence. Intuitively, the integer $\Enclose{m\mapsto n}$ with
$n,m\in \NN$ represents the number $n-m\in \ZZ$. Addition on $\ZZ$ is easily
defined:
\[
 \Enclose{m \mapsto n}+\Enclose{M\mapsto N} = \Enclose{m+M \mapsto n+N}.
\]
\mynote{If integers are interpreted as \emph{translations} on natural numbers, the sum of two integers coincides, where defined, with the composition of the two translations.}
%
It is easily verified that the previous definition is \emph{well-defined},
meaning it does not depend on the choice of $(n,m)$ and $(N,M)$ within their
equivalence class: if $n+m' = n'+m$ and $N+M'=N'+M$, then
$(n+N)+(m'+M')=(n'+N')+(m+M)$.

The addition we have defined on $\ZZ$ is commutative and associative, the number
$\Enclose{0 \mapsto 0}$ is the identity element, and every element has an
inverse, indeed $\Enclose{m \mapsto n} + \Enclose{n \mapsto m} = \Enclose{m+n
\mapsto n+m} = \Enclose{0\mapsto 0}$, that is $-\Enclose{m \mapsto
n}=\Enclose{n\mapsto m}$. Thus, $\ZZ$ is an additive group.

We can easily define an order on $\ZZ$: $\Enclose{m\mapsto n} \ge
\Enclose{M\mapsto N}$ if $n+M\ge N+m$. It is very easy to verify that the order
is compatible with addition, and therefore $\ZZ$ turns out to be a totally
ordered abelian group (definition~\ref{def:gruppo_ordinato}).

We observe that we can identify a natural number $n\in \NN$ with the equivalence
class $\phi(n) = \Enclose{0\mapsto n}\in \ZZ$. The correspondence $\phi\colon
\NN\to\ZZ$ is injective and respects both addition: $\phi(n+m) =
\phi(n)+\phi(m)$ and order: $n\ge m \iff \phi(n)\ge \phi(m)$. Thus, the set
$\phi(\NN)\subset \ZZ$ is an isomorphic copy of the natural numbers and
therefore satisfies the Peano axioms, like $\NN$.
\mynote{
If $n,m\in \NN$, $n\ge m$ it follows that $\Enclose{m \mapsto n} =
\Enclose{n\mapsto 0}-\Enclose{m\mapsto 0} = \phi(n)-\phi(m)$.
}%
It will therefore be convenient to rename $\NN$ by identifying it with
$\phi(\NN)$, thus obtaining $\NN\subset \ZZ$. For any $n\in \ZZ$, it will be
that $n\in \NN$ if $n\ge 0\in \ZZ$, otherwise $-n\in \NN$. Thus, $\ZZ = \NN \cup
(-\NN)$ and $0\in \NN\subset \ZZ$ is the only integer that is its own inverse.

The multiplication operation we defined on $\NN$ can be naturally extended to
all of $\ZZ$. If we want to maintain the distributive property, we must have,
for $n,m\in\NN$
\[
  0 = 0 \cdot n = (m+(-m))\cdot n = m\cdot n + (-m)\cdot n
\]
therefore, by definition, for every $n,m\in \NN$ we set:
\[
  (-m) \cdot n = -(m\cdot n), \qquad  m \cdot (-n) = -(m\cdot n).
\]
In this way, the multiplication $x\cdot y$ is well-defined for every $x,y\in
\ZZ$ and satisfies the distributive property.

With the given definitions, it is tedious but easy to prove that the usual
properties of multiplication hold on $\ZZ$.

\begin{theorem}[Properties of Multiplication]
  Multiplication on $\ZZ$ satisfies the following properties.
  \begin{enumerate}
    \item[1.] Identity and absorbing element: $n\cdot 1 = n$, $n\cdot 0 = 0$,
    \item[2.] Associative property: $(n\cdot m)\cdot k = n \cdot (m\cdot k)$,
    \item[3.] Commutative property: $n\cdot m = m\cdot n$,
    \item[4.] Distributive property: $k\cdot(m+n) = k\cdot m + k\cdot n$. 
  \end{enumerate}
\end{theorem}

Indeed, we discover that $\ZZ$ is an abelian ring with unity
(definition~\ref{def:anello}).

If $n,m\in\ZZ$ with $m\neq 0$ and if there exists $k\in\ZZ$ such that $n=km$, we
will say that $n$ is a \emph{multiple} of $m$ or that $m$ is a \emph{divisor} of
$n$, and we will write:
\[
  \frac{n}{m} = k.  
\]

Multiples of $2$ are called \emph{even} numbers, and integers that are not even
are called \emph{odd} numbers.

\subsection{Fractions and Rational Numbers}

Just as we extended the natural numbers so that the operation of addition has an
inverse operation (subtraction), similarly, we can extend the set of integers so
that multiplication also has an inverse operation (division). The fraction
$\frac p q$ can be defined as a pair of integers $p\in \ZZ$ and $q\in
\NN\setminus \ENCLOSE 0$. The set of rational numbers is obtained by identifying
equivalent fractions, that is:
\[
  \QQ = ((\NN \setminus\ENCLOSE 0)\times \ZZ)/\sim  
\]
where the equivalence relation $\sim$ is defined by setting $(q,p) \sim (q',p')$
when $p\cdot q' = p'\cdot q$. The pair $(q,p)$ represents the fraction $\frac p
q$, and for greater clarity, we will write $\frac p q$ instead of $(q,p)$.

\mynote{The pair $(q,p)$ can also be denoted by the arrow $q\mapsto p$ and represents the dilation of $\ZZ$ that maps the number $q$ to the number $p$. Equivalent fractions represent the same dilation: for example, if I dilate $\ZZ$ by a factor of $\frac 3 2$, the number $2$ will go to the number $3$, and the number $8$ will go to the number $12$. In this sense, the fraction $\frac 3 2$ is considered equivalent to the fraction $\frac{12} 8$. The multiplication of fractions corresponds to the composition of scalings: if I compose the dilation $2\mapsto 3$ with the dilation $3\mapsto 5$, I obtain the dilation $2\mapsto 5$. Indeed, $\frac{3}{2}\cdot \frac{5}{3} = \frac{5}{2}$.}%
Multiplication and addition of rational numbers are defined on the equivalence
classes:
\[
 \Enclose{\frac p q} \cdot \Enclose{\frac{p'}{q'}} 
 = \Enclose{\frac{p\cdot p'}{q\cdot q'}} 
 \qquad 
 \Enclose{\frac{p}{q}} + \Enclose{\frac{p'}{q'}} 
 = \Enclose{\frac{p\cdot q' + p'\cdot q}{q\cdot q'}}.
\]
We can also define an order by setting
\[
 \Enclose{\frac p q} \le \Enclose{\frac{p'}{q'}}
 \quad\iff  \quad p \cdot q' \le p' \cdot q.
\]
It is easily verified that these definitions \emph{pass to the quotient} in the
sense that the product of equivalent fractions is equivalent to the product of
the fractions. Thus, multiplication, addition, and the order relation are
well-defined on $\QQ$.

Fractions of the type $\frac{q}{q}$ are all equivalent to each other and
represent the identity element of multiplication in $\QQ$. Fractions of the type
$\frac{p}{1}$ represent the integers $\ZZ$ within $\QQ$.

With the operations just defined, it can be verified that $\QQ$ is an ordered
and dense field. Moreover, by identifying the rational number $\Enclose{\frac p
1}_\sim$ with the integer $p\in \ZZ$, we can think, by appropriately redefining
$\ZZ$, that $\ZZ \subset \QQ$.

An operation that cannot be defined on $\QQ$ is, for example, the extraction of
the square root. The following theorem tells us, for example, that we cannot
define the square root of $2$ in $\QQ$.

\begin{theorem}[Pythagoras, Irrationality of $\sqrt 2$]
  \mymark{**}%
  \label{th:pitagora}%
  The equation $x^2=2$ has no solutions in $\QQ$.
  \end{theorem}
  %
  \begin{proof}
  \mymark{*}%
    Suppose $x\in \QQ$ is a solution of $x^2=2$. Then it can be written as
  $x=p/q$ with $p\in \ZZ$ and $q\in \NN$, $q\neq 0$. We can also assume that the
  fraction $p/q$ is in its simplest form, meaning that $p$ and $q$ have no
  common factors. Multiplying the equation $(p/q)^2=2$ by $q^2$ yields $p^2 = 2
  q^2$. It follows that $p^2$ is even. But then $p$ must also be even (because
  the square of an odd number is odd). But if $p$ is even, then $p^2$ is a
  multiple of four. But then $2q^2$ is also a multiple of four, and therefore
  $q^2$ is even. Thus, $q$ must also be even. But we assumed that $p$ and $q$
  have no common factors, so this cannot happen.
\end{proof}

The fact that $x^2=2$ has no solution in $\QQ$ allows us to assert that the
order of $\QQ$ is not continuous. Indeed, we can consider the following subsets
of $\QQ$:
\[
 A = \ENCLOSE{x\in \QQ\colon x^2\le 2, x\ge 0}, \qquad
 B = \ENCLOSE{x\in \QQ\colon x^2\ge 2, x\ge 0}. 
\]
Clearly, if $a\in A$ and $b\in B$, we have $a\le b$ because if instead $a>b$, we
would have, by monotonicity, $a^2>b^2$, and thus it could not be $a^2 \le 2$ and
$b^2\ge 2$. Obviously, $0^2 \le 2$ while $2^2\ge 2$, so neither $A$ nor $B$ is
empty. If $\QQ$ were continuous, there should then be a separation element $c$,
that is, $c\in \QQ$ such that $a\le c\le b$ for every $a\in A$ and $b\in B$. We
then ask whether $c^2$ is less than or greater than $2$. If $c^2<2$, we can set
$\eps = 2-c^2 > 0$. There exists a rational number $\delta>0$ such that
$\delta<1$ and $\delta < \frac{\eps}{2c+1}$. Then we would have:
\[
 (c+\delta)^2 
  = c^2+2c\delta + \delta^2 
  < c^2 + 2c\delta + \delta 
  = c^2 + \delta\cdot(2c+1) 
  < c^2 + \eps = 2.
\]
But this means that $c+\delta \in A$, and therefore $c$ cannot be an upper bound
of $A$. If instead $c^2>2$, since $1\in A$, it must also be $c\ge 1>0$. Thus, we
can take $\delta = \frac{c^2-2}{2c}>0$, and we would have
\[
 (c-\delta)^2 
 = c^2 - 2c\delta + \delta^2
 > c^2 - 2c\delta = 2
\]
from which it follows that $c-\delta \in B$, and therefore $c$ cannot be a lower
bound of $B$. We must therefore conclude that it must be $c^2=2$. But, by
theorem~\ref{th:pitagora}, this is not possible. This means that $A$ and $B$ do
not have a separation element in $\QQ$.